% Section 9: Cosmology
\section{Cosmology}
\label{sec:cosmology}

\subsection{Torsion-Driven Inflation}

\subsubsection{The Early Universe}

In the very early universe ($t < 10^{-43}$ s), the spectral dimension was $d_s \approx 10$, and spacetime had a highly fractal structure. The torsion field was correspondingly strong:

\begin{equation}
    \tau_{early} \sim \tau_0 \left(\frac{E}{M_P}\right)^{d_s - 4} \sim \tau_0
\end{equation}

\subsubsection{Inflationary Dynamics}

The torsion field acts as an effective inflaton, driving exponential expansion:

\begin{equation}
    H^2 = \frac{8\pi G}{3}\left(\rho_{matter} + \rho_{torsion}\right)
\end{equation}

where the torsion energy density is:
\begin{equation}
    \rho_{torsion} = \frac{1}{2}\tau^2 + \frac{1}{12}\tau^4
\end{equation}

\subsubsection{Slow-Roll Conditions}

For successful inflation, the slow-roll parameters are:
\begin{align}
    \epsilon &= -\frac{\dot{H}}{H^2} = \frac{1}{16\pi G}\left(\frac{V'}{V}\right)^2 \ll 1 \\
    \eta &= \frac{1}{8\pi G}\frac{V''}{V} \ll 1
\end{align}

The torsion potential naturally satisfies these conditions when $\tau \sim \tau_0$.

\subsubsection{Number of E-folds}

The number of e-folds during inflation is:
\begin{equation}
    N = \int_{t_i}^{t_f} H dt \approx 60
\end{equation}

This naturally explains the flatness and horizon problems.

\subsection{Spectral Dimension Evolution in the Early Universe}

\subsubsection{The Transition $d_s: 10 \to 4$}

As the universe expands and cools, the spectral dimension flows from 10 to 4:

\begin{equation}
    d_s(T) = 4 + \frac{6}{1 + (T/T_c)^2}
\end{equation}

where $T_c \sim 10^{16}$ GeV is the GUT scale.

\subsubsection{Timeline of Dimensionality Changes}

\begin{center}
\begin{tabular}{llll}
\toprule
\textbf{Epoch} & \textbf{Time} & \textbf{$d_s$} & \textbf{Physics} \\
\midrule
Planck & $t < 10^{-43}$ s & 10 & Quantum gravity \\
String & $10^{-43}$--$10^{-40}$ s & 9--10 & String excitations \\
GUT & $10^{-40}$--$10^{-36}$ s & 7--9 & Symmetry breaking \\
Inflation & $10^{-36}$--$10^{-32}$ s & 5--7 & Rapid expansion \\
Radiation & $10^{-32}$--$10^{4}$ s & 4 & Standard cosmology \\
Matter & $t > 10^{4}$ s & 4 & Structure formation \\
\bottomrule
\end{tabular}
\end{center}

\subsubsection{Effect on Cosmic Evolution}

The changing dimensionality affects:
\begin{enumerate}
    \item \textbf{Expansion rate}: $H \propto d_s^{1/2}$
    \item \textbf{Particle production}: $\Gamma \propto T^{d_s}$
    \item \textbf{Entropy density}: $s \propto T^{d_s-1}$
\end{enumerate}

\subsection{Geometric Dark Energy}

\subsubsection{The Cosmological Constant Problem}

Standard quantum field theory predicts a cosmological constant:
\begin{equation}
    \Lambda_{QFT} \sim M_P^4 \sim 10^{76} \text{ GeV}^4
\end{equation}

while observation gives:
\begin{equation}
    \Lambda_{obs} \sim 10^{-47} \text{ GeV}^4
\end{equation}

This 123-order-of-magnitude discrepancy is the cosmological constant problem.

\subsubsection{Fractal Suppression Mechanism}

In our theory, the fractal structure at small scales suppresses the vacuum energy:

\begin{equation}
    \Lambda_{eff} = \Lambda_{bare} \cdot \left(\frac{\ell_P}{L}\right)^{d_s - 4}
\end{equation}

For $d_s \approx 10$ at the Planck scale:
\begin{equation}
    \Lambda_{eff} \sim \Lambda_{bare} \cdot 10^{-123}
\end{equation}

\subsubsection{Dark Energy as Geometric Effect}

The observed dark energy is not a cosmological constant but a geometric effect of the fractal structure:

\begin{equation}
    \rho_{DE} = \tau_0^2 \cdot f(z)
\end{equation}

where $f(z)$ is a slowly varying function of redshift.

This gives an effective equation of state:
\begin{equation}
    w = \frac{p}{\rho} \approx -1 + \mathcal{O}(\tau_0^2)
\end{equation}

consistent with observations.

\subsection{Primordial Gravitational Waves}

\subsubsection{Generation Mechanism}

Inflation generates primordial gravitational waves with spectrum:
\begin{equation}
    \Omega_{GW}(f) = \Omega_{GW}^{(0)} \left(\frac{f}{f_*}
ight)^{n_T}
\end{equation}

where the tensor spectral index is:
\begin{equation}
    n_T = -2\epsilon \approx -\frac{r}{8}
\end{equation}

\subsubsection{Six Polarization Modes in the Early Universe}

During the high-$d_s$ phase, all six polarization modes were active:
\begin{equation}
    h_{total} = h_+ + h_\times + h_x + h_y + h_b + h_l
\end{equation}

The additional modes contribute to the stochastic gravitational wave background.

\subsubsection{Detection Prospects}

Future detectors could observe:
\begin{itemize}
    \item \textbf{CMB B-modes}: Primordial tensor perturbations
    \item \textbf{PTA}: Stochastic background from supermassive black holes
    \item \textbf{LISA}: Primordial phase transitions at GUT scale
    \item \textbf{Atom interferometers}: High-frequency primordial signals
\end{itemize}

\subsection{Structure Formation}

\subsubsection{Power Spectrum}

The primordial power spectrum receives corrections from the fractal structure:
\begin{equation}
    P(k) = P_{\Lambda CDM}(k) \cdot (1 + \alpha \tau_0^2 \ln(k/k_*))
\end{equation}

where $\alpha$ is a numerical coefficient.

\subsubsection{Dark Matter}

Quantum black hole remnants contribute to dark matter:
\begin{equation}
    \Omega_{rem} = \frac{n_{rem} M_{rem}}{\rho_c}
\end{equation}

The abundance depends on the primordial black hole formation rate.

\subsection{Summary}

Cosmology in our framework features:
\begin{enumerate}
    \item \textbf{Torsion-driven inflation}: Natural slow-roll without fine-tuning
    \item \textbf{Dynamic spectral dimension}: $d_s: 10 \to 4$ through cosmic history
    \item \textbf{Geometric dark energy}: Resolution of cosmological constant problem
    \item \textbf{Primordial gravitational waves}: Six polarization modes in early universe
    \item \textbf{Remnant dark matter}: Evaporated black holes as dark matter candidates
\end{enumerate}

These predictions are testable with current and future cosmological observations.
